\documentclass[11pt,a4paper,twoside]{book}

\usepackage[utf8]{inputenc}
\usepackage[T1]{fontenc}
\usepackage[french]{babel}
\usepackage{lmodern}
\usepackage[a4paper,margin=2.2cm]{geometry}
\usepackage{amsmath,amssymb}
\usepackage{array,booktabs,longtable}
\usepackage{enumitem}
\usepackage[dvipsnames,svgnames]{xcolor}
\usepackage[most]{tcolorbox}
\usepackage{fancyhdr}
\usepackage[hidelinks]{hyperref}
\setlength{\headheight}{13pt}

\definecolor{capbrown}{HTML}{B45309}
\definecolor{caplight}{HTML}{FFF7ED}
\definecolor{capgold}{HTML}{FDE68A}
\definecolor{capblue}{HTML}{0369A1}
\definecolor{capbluebg}{HTML}{F0F9FF}
\definecolor{capgreen}{HTML}{15803D}
\definecolor{capgreenbg}{HTML}{F0FDF4}
\definecolor{capred}{HTML}{B91C1C}
\definecolor{capredbg}{HTML}{FEF2F2}
\definecolor{capgray}{HTML}{475569}

\newtcolorbox{definitionbox}[1][]{%
  breakable,
  colback=caplight,
  colframe=capbrown,
  fonttitle=\bfseries,
  title={Définition\ifx\relax#1\relax\else\space--- #1\fi}
}

\newtcolorbox{methodebox}[1][]{%
  breakable,
  colback=capbluebg,
  colframe=capblue,
  fonttitle=\bfseries,
  title={Méthode\ifx\relax#1\relax\else\space--- #1\fi}
}

\newtcolorbox{exemplebox}[1][]{%
  breakable,
  colback=capgreenbg,
  colframe=capgreen,
  fonttitle=\bfseries,
  title={Exemple\ifx\relax#1\relax\else\space--- #1\fi}
}

\newtcolorbox{attentionbox}[1][]{%
  breakable,
  colback=capredbg,
  colframe=capred,
  fonttitle=\bfseries,
  title={Attention\ifx\relax#1\relax\else\space--- #1\fi}
}

\newtcolorbox{retenirbox}{%
  breakable,
  colback=white,
  colframe=capbrown,
  coltitle=white,
  colbacktitle=capbrown,
  fonttitle=\bfseries,
  title={À retenir},
  boxrule=1.2pt
}

\setlist[itemize]{topsep=4pt,itemsep=3pt}
\setlist[enumerate]{topsep=4pt,itemsep=3pt}

\pagestyle{fancy}
\fancyhf{}
\fancyhead[LE]{\small\sffamily\leftmark}
\fancyhead[RO]{\small\sffamily\rightmark}
\fancyfoot[C]{\thepage}
\renewcommand{\headrulewidth}{0.4pt}

\begin{document}

\begin{titlepage}
\centering
\vspace*{2.5cm}
{\Huge\bfseries Cours de Mathématiques\par}
\vspace{0.4cm}
{\LARGE CAP\par}
\vspace{0.4cm}
{\Large Groupement 1\par}
\vspace{1.5cm}
{\large MIT, Ébéniste, Signalétique et Décors Graphiques\par}
\vspace{2cm}
\begin{tcolorbox}[colback=caplight,colframe=capbrown,width=.85\textwidth]
Ce livre regroupe les cours de mathématiques CAP présents dans le site
\texttt{maths-sciences-lp.github.io}. Il est structuré en 7 chapitres :
statistiques, probabilités, proportionnalité, équations, fonctions,
géométrie et calculs numériques.
\end{tcolorbox}
\vfill
{\large M. Azzouz\par}
{\large Lycée Eugène Hénaff\par}
\vspace{0.5cm}
{\large \today\par}
\end{titlepage}

\frontmatter
\tableofcontents

\mainmatter

\chapter{Statistiques à une variable}

\section{Pourquoi les statistiques ?}
\begin{exemplebox}[Situation professionnelle]
Un artisan menuisier mesure la longueur de pièces de bois destinées à des placards sur mesure.
Il cherche à savoir quelle valeur revient le plus souvent et quelle est la longueur moyenne.
Les statistiques permettent de recueillir, organiser, représenter puis interpréter ces données.
\end{exemplebox}

\section{Vocabulaire}
\begin{definitionbox}
Une \textbf{série statistique} est un ensemble de données recueillies lors d'une étude.
\begin{itemize}
\item La \textbf{population} est l'ensemble des individus étudiés.
\item Le \textbf{caractère} est la propriété observée.
\item Les \textbf{valeurs} sont les résultats possibles du caractère.
\end{itemize}
\end{definitionbox}

\section{Effectifs et fréquences}
\begin{definitionbox}
L'\textbf{effectif} d'une valeur est le nombre de fois où cette valeur apparaît.
L'\textbf{effectif total} est noté \(N\).

La \textbf{fréquence} d'une valeur est donnée par :
\[
\text{fréquence}=\frac{\text{effectif de la valeur}}{\text{effectif total}}
\]
\end{definitionbox}

\begin{methodebox}[Construire un tableau statistique]
\begin{enumerate}
\item Lister les valeurs différentes.
\item Compter les effectifs.
\item Calculer les fréquences.
\item Vérifier que la somme des effectifs vaut \(N\) et que la somme des fréquences vaut \(1\).
\end{enumerate}
\end{methodebox}

\begin{exemplebox}
Pour les températures \(55,58,55,60,58,55,62,58\), on a \(N=8\).
\[
\begin{array}{c|cccc}
\text{Température} & 55 & 58 & 60 & 62\\ \hline
\text{Effectif} & 3 & 3 & 1 & 1
\end{array}
\]
La fréquence de \(55\) est \(\frac{3}{8}=0{,}375=37{,}5\%\).
\end{exemplebox}

\section{Représentations graphiques}
\begin{definitionbox}
Un \textbf{diagramme en bâtons} représente des données discrètes par des segments verticaux.

Un \textbf{diagramme circulaire} représente chaque valeur par un secteur d'angle
proportionnel à son effectif ou à sa fréquence.
\end{definitionbox}

\begin{methodebox}[Angle d'un secteur]
\[
\text{angle}=\text{fréquence}\times 360^\circ
\qquad\text{ou}\qquad
\text{angle}=\frac{\text{effectif}}{N}\times 360^\circ
\]
\end{methodebox}

\begin{attentionbox}
Dans un diagramme circulaire, la somme des angles doit toujours faire \(360^\circ\).
\end{attentionbox}

\section{Moyenne}
\begin{definitionbox}
La moyenne d'une série statistique est
\[
\bar{x}=\frac{\text{somme des valeurs}}{\text{effectif total}}
\]
\end{definitionbox}

\begin{exemplebox}
Pour les notes \(8,10,12,14,16\),
\[
\bar{x}=\frac{8+10+12+14+16}{5}=12
\]
\end{exemplebox}

\section{Lire et interpréter un graphique}
\begin{methodebox}
\begin{enumerate}
\item Lire le titre et les axes.
\item Identifier l'unité utilisée.
\item Relever les valeurs ou les comparer.
\item Conclure avec une phrase rédigée.
\end{enumerate}
\end{methodebox}

\begin{retenirbox}
\begin{itemize}
\item Une série statistique se décrit avec des effectifs, des fréquences et des indicateurs.
\item Les fréquences totalisent \(1\), soit \(100\%\).
\item On choisit une représentation adaptée à la nature des données.
\item La moyenne donne une valeur globale de la série.
\end{itemize}
\end{retenirbox}

\chapter{Probabilités}

\section{Expérience aléatoire et issues}
\begin{definitionbox}
Une \textbf{expérience aléatoire} est une expérience dont on ne peut pas prévoir à l'avance
le résultat exact, mais dont on connaît les résultats possibles.

Les \textbf{issues} sont tous les résultats possibles de cette expérience.
\end{definitionbox}

\begin{exemplebox}
Lancer un dé à 6 faces : les issues sont \(1,2,3,4,5,6\).
\end{exemplebox}

\section{Événements}
\begin{definitionbox}
Un \textbf{événement} est une condition portant sur le résultat d'une expérience aléatoire.
\end{definitionbox}

\begin{exemplebox}
Avec un dé :
\begin{itemize}
\item obtenir un \(3\) ;
\item obtenir un nombre pair ;
\item obtenir un nombre inférieur à \(5\).
\end{itemize}
\end{exemplebox}

\section{Probabilité d'un événement}
\begin{definitionbox}
La \textbf{probabilité} d'un événement est un nombre compris entre \(0\) et \(1\).
\end{definitionbox}

\begin{methodebox}[Cas équiprobable]
Si toutes les issues ont la même chance de se produire :
\[
P(\text{événement})=\frac{\text{nombre d'issues favorables}}{\text{nombre total d'issues}}
\]
\end{methodebox}

\begin{exemplebox}
Avec un dé équilibré, la probabilité d'obtenir un nombre pair est :
\[
P(\text{pair})=\frac{3}{6}=\frac{1}{2}=0{,}5=50\%
\]
\end{exemplebox}

\begin{attentionbox}
Une probabilité ne peut jamais être négative ni supérieure à \(1\).
\end{attentionbox}

\section{Événement contraire}
\begin{definitionbox}
L'événement contraire de \(A\), noté \(\overline{A}\), est l'événement qui se réalise
quand \(A\) ne se réalise pas.
\end{definitionbox}

\begin{methodebox}
\[
P(\overline{A})=1-P(A)
\]
\end{methodebox}

\section{Fluctuation et stabilisation}
\begin{definitionbox}
Quand on répète une expérience un petit nombre de fois, la fréquence observée fluctue.
Quand le nombre d'essais augmente, elle se rapproche de la probabilité théorique.
\end{definitionbox}

\begin{retenirbox}
\begin{itemize}
\item Une expérience aléatoire possède des issues possibles.
\item Un événement peut être certain, impossible ou simplement probable.
\item En situation équiprobable, on utilise le rapport \(\frac{\text{favorables}}{\text{total}}\).
\item La fréquence observée se stabilise autour de la probabilité théorique quand le nombre d'essais augmente.
\end{itemize}
\end{retenirbox}

\chapter{Proportionnalité et pourcentages}

\section{Suites proportionnelles}
\begin{definitionbox}
Deux suites sont \textbf{proportionnelles} quand on passe de l'une à l'autre
en multipliant toujours par le même nombre, appelé \textbf{coefficient de proportionnalité}.
\end{definitionbox}

\begin{methodebox}
Pour vérifier une proportionnalité, on compare les quotients correspondants.
S'ils sont tous égaux, la situation est proportionnelle.
\end{methodebox}

\begin{exemplebox}
Des tasseaux coûtent \(2{,}50\) euros pièce.
\[
\frac{2{,}50}{1}=\frac{5}{2}=\frac{12{,}50}{5}=2{,}50
\]
Le prix est proportionnel au nombre de pièces.
\end{exemplebox}

\section{Quatrième proportionnelle}
\begin{definitionbox}
Dans un tableau de proportionnalité où trois valeurs sont connues,
la quatrième se calcule par produit en croix.
\end{definitionbox}

\begin{methodebox}[Règle de trois]
Si \(a\) correspond à \(b\), alors la valeur \(x\) correspondant à \(c\) vérifie :
\[
x=\frac{b\times c}{a}
\]
\end{methodebox}

\begin{exemplebox}
Si \(3\) m de tuyau coûtent \(18\) euros, alors \(7\) m coûtent :
\[
x=\frac{18\times 7}{3}=42
\]
\end{exemplebox}

\section{Pourcentages}
\begin{definitionbox}
Un pourcentage est une proportion pour \(100\).
\[
t\%= \frac{t}{100}
\]
\end{definitionbox}

\begin{methodebox}
Pour calculer \(t\%\) d'une quantité \(Q\) :
\[
Q\times \frac{t}{100}
\]
\end{methodebox}

\begin{exemplebox}
La TVA de \(10\%\) sur \(168\) euros vaut :
\[
168\times 0{,}10=16{,}80
\]
Le prix TTC vaut donc \(184{,}80\) euros.
\end{exemplebox}

\section{Coefficient multiplicateur}
\begin{definitionbox}
\begin{itemize}
\item Pour une augmentation de \(t\%\), on multiplie par \(1+\frac{t}{100}\).
\item Pour une réduction de \(t\%\), on multiplie par \(1-\frac{t}{100}\).
\end{itemize}
\end{definitionbox}

\begin{exemplebox}
Après une remise de \(15\%\), le coefficient multiplicateur est \(0{,}85\).
\end{exemplebox}

\section{Échelles}
\begin{definitionbox}
Une échelle traduit le rapport entre une longueur sur un plan et la longueur réelle.
\end{definitionbox}

\begin{exemplebox}
À l'échelle \(1/50\), \(4\) cm sur le plan représentent \(4\times 50=200\) cm, soit \(2\) m en réalité.
\end{exemplebox}

\begin{retenirbox}
\begin{itemize}
\item La proportionnalité repose sur un coefficient constant.
\item La règle de trois ne s'applique qu'en situation proportionnelle.
\item Un pourcentage s'exprime facilement en décimal.
\item Les augmentations et réductions se traitent efficacement avec un coefficient multiplicateur.
\end{itemize}
\end{retenirbox}

\chapter{Équations du premier degré}

\section{Qu'est-ce qu'une équation ?}
\begin{definitionbox}
Une \textbf{équation} est une égalité contenant une inconnue, souvent notée \(x\).
Résoudre une équation, c'est trouver la valeur de \(x\) qui rend l'égalité vraie.
\end{definitionbox}

\begin{exemplebox}
Dans \(2x+5=11\), la solution est \(x=3\).
\end{exemplebox}

\section{Règles de transformation}
\begin{methodebox}
On peut transformer une équation sans changer sa solution en :
\begin{itemize}
\item ajoutant ou soustrayant le même nombre des deux côtés ;
\item multipliant ou divisant les deux côtés par un même nombre non nul.
\end{itemize}
\end{methodebox}

\section{Résoudre \(ax+b=c\)}
\begin{methodebox}
\begin{enumerate}
\item Isoler le terme en \(x\) : \(ax=c-b\).
\item Diviser par \(a\) : \(x=\frac{c-b}{a}\).
\item Vérifier la solution dans l'équation de départ.
\end{enumerate}
\end{methodebox}

\begin{exemplebox}
Résoudre \(5x+3=28\) :
\[
5x=25 \qquad\Rightarrow\qquad x=5
\]
\end{exemplebox}

\section{Résoudre \(ax+b=cx+d\)}
\begin{methodebox}
\begin{enumerate}
\item Regrouper les termes en \(x\) d'un côté.
\item Regrouper les nombres de l'autre côté.
\item Résoudre l'équation obtenue.
\end{enumerate}
\end{methodebox}

\begin{exemplebox}
\[
7x+4=3x+20
\Rightarrow 4x=16
\Rightarrow x=4
\]
\end{exemplebox}

\section{Modéliser un problème}
\begin{exemplebox}[Situation professionnelle]
Une planche de \(2{,}40\) m est découpée en \(3\) morceaux égaux avec \(2\) traits de scie de \(3\) mm.
Si \(x\) désigne la longueur d'un morceau :
\[
3x+2\times 0{,}003=2{,}40
\]
On obtient :
\[
3x=2{,}394 \qquad\Rightarrow\qquad x=0{,}798\text{ m}
\]
\end{exemplebox}

\begin{attentionbox}
Toujours vérifier le résultat et garder une cohérence d'unités.
\end{attentionbox}

\begin{retenirbox}
\begin{itemize}
\item Une équation sert à déterminer une valeur inconnue.
\item Les transformations doivent être faites des deux côtés.
\item Une résolution correcte se termine par une vérification.
\item La mise en équation est essentielle dans les problèmes concrets.
\end{itemize}
\end{retenirbox}

\chapter{Fonctions}

\section{Notion de fonction}
\begin{definitionbox}
Une \textbf{fonction} associe à chaque nombre \(x\) une unique valeur notée \(f(x)\).
\end{definitionbox}

\begin{definitionbox}[Image et antécédent]
\begin{itemize}
\item \(f(x)\) est l'\textbf{image} de \(x\).
\item \(x\) est un \textbf{antécédent} de \(f(x)\).
\end{itemize}
\end{definitionbox}

\section{Tableau de valeurs}
\begin{methodebox}
Pour compléter un tableau de valeurs :
\begin{enumerate}
\item repérer la formule ;
\item remplacer \(x\) par chaque valeur ;
\item calculer l'image correspondante.
\end{enumerate}
\end{methodebox}

\begin{exemplebox}
Si \(C(x)=0{,}05x+4\), alors :
\[
C(200)=0{,}05\times 200+4=14
\]
\end{exemplebox}

\section{Représentation graphique}
\begin{definitionbox}
La courbe représentative d'une fonction est l'ensemble des points
de coordonnées \((x\,;\,f(x))\) dans un repère.
\end{definitionbox}

\begin{methodebox}
Pour tracer une courbe :
\begin{enumerate}
\item préparer un tableau de valeurs ;
\item choisir des axes gradués adaptés ;
\item placer les points ;
\item relier correctement selon la situation.
\end{enumerate}
\end{methodebox}

\section{Variations}
\begin{definitionbox}
\begin{itemize}
\item Une fonction est \textbf{croissante} quand \(x\) augmente et \(f(x)\) augmente.
\item Elle est \textbf{décroissante} quand \(x\) augmente et \(f(x)\) diminue.
\end{itemize}
\end{definitionbox}

\section{Fonction linéaire}
\begin{definitionbox}
Une fonction linéaire est de la forme :
\[
f(x)=ax
\]
où \(a\) est un nombre réel.
\end{definitionbox}

\begin{exemplebox}
Le coût de planches à \(12\) euros le mètre se modélise par :
\[
C(x)=12x
\]
\end{exemplebox}

\begin{attentionbox}
Chaque \(x\) a une seule image, mais une même image peut avoir plusieurs antécédents selon la fonction étudiée.
\end{attentionbox}

\begin{retenirbox}
\begin{itemize}
\item Une fonction relie une grandeur d'entrée à une grandeur de sortie.
\item Les tableaux et les graphiques sont deux représentations complémentaires.
\item Les variations aident à interpréter l'évolution d'une grandeur.
\item La fonction linéaire modélise les situations proportionnelles.
\end{itemize}
\end{retenirbox}

\chapter{Géométrie}

\section{Figures planes, périmètres et aires}
\begin{definitionbox}
Formules usuelles :
\[
\mathcal{P}_{\text{rectangle}}=2(L+\ell), \qquad
\mathcal{A}_{\text{rectangle}}=L\times \ell
\]
\[
\mathcal{P}_{\text{cercle}}=2\pi r, \qquad
\mathcal{A}_{\text{disque}}=\pi r^2
\]
\end{definitionbox}

\begin{exemplebox}
Un panneau de \(1{,}20\) m sur \(0{,}80\) m a pour aire :
\[
1{,}20\times 0{,}80=0{,}96\text{ m}^2
\]
\end{exemplebox}

\section{Volumes}
\begin{definitionbox}
Formules usuelles :
\[
V_{\text{pavé}}=L\times \ell \times h,\qquad
V_{\text{cylindre}}=\pi r^2 h,\qquad
V_{\text{boule}}=\frac{4}{3}\pi r^3
\]
\end{definitionbox}

\begin{exemplebox}
Pour un cylindre de rayon \(0{,}25\) m et de hauteur \(1{,}20\) m :
\[
V=\pi \times 0{,}25^2 \times 1{,}20 \approx 0{,}236\text{ m}^3
\]
soit environ \(236\) L.
\end{exemplebox}

\section{Théorème de Pythagore}
\begin{definitionbox}
Dans un triangle rectangle :
\[
\text{hypoténuse}^2=\text{côté 1}^2+\text{côté 2}^2
\]
\end{definitionbox}

\begin{exemplebox}
Pour un rectangle de côtés \(60\) cm et \(80\) cm :
\[
d^2=60^2+80^2=10000
\Rightarrow d=100\text{ cm}
\]
\end{exemplebox}

\section{Théorème de Thalès}
\begin{definitionbox}
Le théorème de Thalès permet de comparer des longueurs dans des configurations de droites parallèles.
\end{definitionbox}

\begin{methodebox}
Avant d'utiliser Thalès :
\begin{enumerate}
\item repérer les droites parallèles ;
\item écrire les longueurs dans le bon ordre ;
\item vérifier la cohérence des rapports.
\end{enumerate}
\end{methodebox}

\section{Symétrie}
\begin{definitionbox}
Deux figures sont symétriques si l'une peut se superposer à l'autre après réflexion par rapport à un axe.
\end{definitionbox}

\section{Conversions d'unités}
\begin{attentionbox}
\begin{itemize}
\item \(1\text{ m}^2 = 10\,000\text{ cm}^2\)
\item \(1\text{ dm}^3 = 1\text{ L}\)
\item \(1\text{ m}^3 = 1\,000\text{ L}\)
\end{itemize}
\end{attentionbox}

\begin{retenirbox}
\begin{itemize}
\item Les calculs de périmètre, d'aire et de volume doivent toujours être accompagnés d'unités.
\item Pythagore s'utilise dans les triangles rectangles.
\item Thalès s'utilise dans des configurations avec parallèles.
\item Les conversions d'unités sont une source fréquente d'erreurs.
\end{itemize}
\end{retenirbox}

\chapter{Calculs numériques}

\section{Fractions}
\begin{definitionbox}
Une fraction est un quotient de deux entiers :
\[
\frac{a}{b}\qquad \text{avec } b\neq 0
\]
\end{definitionbox}

\begin{methodebox}
\begin{itemize}
\item Même dénominateur : \(\frac{a}{b}+\frac{c}{b}=\frac{a+c}{b}\)
\item Multiplication : \(\frac{a}{b}\times\frac{c}{d}=\frac{ac}{bd}\)
\item Division : \(\frac{a}{b}\div\frac{c}{d}=\frac{a}{b}\times\frac{d}{c}\)
\end{itemize}
\end{methodebox}

\section{Nombres décimaux et priorités}
\begin{definitionbox}
Priorités de calcul :
\begin{enumerate}
\item parenthèses ;
\item puissances ;
\item multiplications et divisions ;
\item additions et soustractions.
\end{enumerate}
\end{definitionbox}

\begin{exemplebox}
\[
3+4\times 5-2=3+20-2=21
\]
\end{exemplebox}

\section{Puissances de 10 et écriture scientifique}
\begin{definitionbox}
L'écriture scientifique d'un nombre est de la forme :
\[
a\times 10^n
\qquad \text{avec} \qquad 1\le a < 10
\]
\end{definitionbox}

\begin{exemplebox}
\[
45000=4{,}5\times 10^4
\qquad\text{et}\qquad
0{,}0032=3{,}2\times 10^{-3}
\]
\end{exemplebox}

\section{Ordres de grandeur et arrondis}
\begin{definitionbox}
Un ordre de grandeur permet d'estimer rapidement un résultat.
Un arrondi donne une valeur approchée à l'unité, au dixième, au centième, etc.
\end{definitionbox}

\section{Conversions d'unités}
\begin{methodebox}
\begin{itemize}
\item Longueurs : \(1\text{ m}=100\text{ cm}=1000\text{ mm}\)
\item Masses : \(1\text{ kg}=1000\text{ g}\)
\item Capacités : \(1\text{ L}=1\text{ dm}^3\)
\end{itemize}
\end{methodebox}

\begin{attentionbox}
Toujours convertir toutes les données dans la même unité avant de calculer.
\end{attentionbox}

\begin{retenirbox}
\begin{itemize}
\item Les fractions et les décimaux doivent être manipulés avec méthode.
\item Les priorités de calcul évitent les erreurs.
\item Les puissances de \(10\) simplifient l'écriture et l'estimation.
\item Une bonne maîtrise des conversions d'unités est indispensable dans les métiers préparés.
\end{itemize}
\end{retenirbox}

\backmatter

\chapter*{Source}
Ce livre a été rédigé à partir des leçons CAP présentes dans :
\begin{center}
\texttt{maths/cap/ch01..ch07/lecon.html}
\end{center}

\end{document}
